\section{연구 방법 및 절차}

\subsection{모델의 가정}

다리나 체인링, 자전거 등이 크게 모양이 역동적으로 변하는 대상은 아니므로 그런 대상들이 모두 강체라고 가정한다.

사람 다리에서의 가정

사람의 다리와 페달은 얇은 질량분포가 균일한 세 막대가 자유로운 관절로 연결되어 있는 것이라고 가정한다. 실제로는 관절이 그렇게 단순하지는 않겠으나 우리 연구가 보고자 하는 것은 체인링의 모양이 바뀔 때 사람의 부담이 어떻게 되느냐이다. 따라서 관절이 어떻게 작동하느냐는 중요도가 떨어지는 요소로 근사적으로 보았다.

사람이 근육으로 내는 힘은 고관절에서만 난다고 가정한다. 실제 자전거를 탈 때 허벅지 근육을 사용하여 고관절을 돌리는 것이 앞으로 진행하는 것의 주가 되기 때문이다.

한 번에 한 다리만 페달과 상호작용한다고 가정한다. 자전거를 탈 때 페달을 다리로 당길 수는 없다. 페달을 굴러 앞으로 가기 위해서는 진행 방향 쪽에 있는 페달을 미는 방법밖에는 없고 진행 방향의 뒷쪽에 있는 페달을 당겨서 앞으로 가는 힘을 제공할 수는 없다. 즉 뒷쪽에 있는 페달로 자전거가 앞으로 가는 것에 기여할 수 없다. 따라서 가장 효율적으로 자전거를 타는 방법은 한 쪽 발은 아예 기여하지 않는 것이다. 우리는 사람이 자전거를 어떻게 타는지에 관심이 없으므로 가장 효율적인 방식으로 자전거를 탄다고 가정하여 한 번에 한 다리만 페달과 상호작용한다고 가정하겠다.

체인링, 스프로켓 시스템에서의 가정

현재 분석하고자 하는 것은 체인링의 형상에 따라 얼마나 피로도가 주는지이므로 체인링이 크거나 작아짐에 의한 효과는 무시하고자 한다. 즉 스프로켓과 체인링의 크기가 같다고 가정을 한다. 따라서 윗 체인은 체인링이 원형이 아님에도 불구하고 항상 수평하다고 가정한다.

체인링이 타원형이면 만약 윗 체인과 아래 체인이 모두 팽팽하다면 체인의 길이가 일정할 수가 없을 것이다. 다만 실제 자전거에서도 기어를 바꿀 때 체인을 팽팽하게 유지시키기 위한 장치인 뒷변속기가 있다. 이와 같은 장치가 충분히 잘 작동하여 체인이 항상 팽팽하게 유지된다고 가정할 수 있다. 또한 뒷변속기는 일반적으로 스프로켓에 위치하므로 체인링과 접하는 체인은 항상 수평하다고 가정한다.

자전거의 움직임에서의 가정

자전거는 뒤로 일정한 외력을 받고 등속으로 움직인다고 가정한다. 실제로 일상적인 상황에서 자전거를 탈 때 공기저항이나 중력에 의해 거의 일정한 외력이 작용할 것이며 거의 등속으로 운동할 것이다.

\subsection{물리적인 식들의 유도}

연구의 목적은 어떤 체인링 형상이 사람의 부담을 가장 줄일 수 있는지이므로 주어진 체인링 형상에서 사람의 부담을 계산할 수 있어야 한다.

따라서 연구의 흐름을 자전거가 일정한 저항력에서 등속으로 움직이기 위해 사람이 페달에 가해야 하는 토크를 계산하고 사람의 다리 모델을 통해 고관절이 가해야 하는 토크를 계산하여 이를 바탕으로 부담을 정의하는 방향으로 잡았다.

\begin{table}[!htb]
\centering
\small
\begin{tabular}{@{}clc@{}}
\toprule
기호 & 의미 & 단위\\
\midrule
\(t\) & 시간 & s\\
\(\theta\) & 크랭크의 회전각 & rad\\
\(\rho(\theta)\) & 크랭크축에서 윗 체인까지의 수직거리 & m\\
\(\rho'(\theta)\) & 크랭크 각도에 따른 \(\rho\)의 변화율 & m/rad\\
\(v_0\) & 자전거의 속력 & m/s\\
\(f\) & 자전거의 운동을 방해하는 저항력 & N\\
\(r,\ R\) & 뒷스프로켓, 바퀴의 반지름 & m\\
\(I_c\) & 크랭크와 체인링의 회전관성 & kg\(\cdot\)m\(^2\)\\
\(L_1,\ L_2\) & 허벅지와 정강이의 길이 & m\\
\(\ell,\ d\) & 크랭크 길이와 고관절--크랭크축 사이의 거리 & m\\
\(m_1,\ m_2\) & 허벅지와 정강이 및 발의 질량 & kg\\
\(\phi\) & 고관절에서 허벅지가 이루는 각 & rad\\
\(k\) & 크랭크 각도에 대한 \(\phi\)의 변화율 & --\\
\(J,\ J'\) & 다리의 등가 회전관성과 그 각도에 대한 변화율 & kg\(\cdot\)m\(^2\), kg\(\cdot\)m\(^2\)/rad\\
\(\tau_p,\ \tau_h\) & 페달이 크랭크에 가하는 토크, 고관절 토크 & N\(\cdot\)m\\
\(\tau_{\mathrm{iso}}\) & 각속도가 0일 때의 최대 고관절 토크 & N\(\cdot\)m\\
\(\omega_{\max}\) & 최대 고관절 토크가 0이 되는 각속도 & rad/s\\
\(u\) & 최대 토크에 대한 순간 고관절 토크의 비율 & --\\
\bottomrule
\end{tabular}
\caption{계산에 사용한 물리량}
\label{tab:physical-quantities}
\end{table}

먼저 자전거가 등속으로 움직일 때 페달에 가해져야 하는 토크를 구하고자 한다. 뒷바퀴가 등속으로 돌아가므로 뒷스프로켓에 걸린 체인의 속력도 일정하다. 체인은 미끄러지지 않으므로 체인링에서의 체인 속력은 크랭크의 각속도와 \(\rho\)의 곱으로 나타낼 수 있다. 따라서 자전거가 같은 속력으로 움직이더라도 \(\rho\)가 변하면 크랭크의 각속도는 변한다.

\begin{equation}
\dot{\theta}=\frac{rv_0}{R\rho},\qquad
\ddot{\theta}=-\frac{r^2v_0^2}{R^2\rho^3}\rho'.
\label{eq:me-motion}
\end{equation}

즉 \(\rho\)가 큰 곳에서는 크랭크가 느리게 돌고, 작은 곳에서는 빠르게 돈다. 체인링이 원형이라면 \(\rho\)가 일정하므로 각속도도 일정하다. 하지만 비원형 체인링에서는 크랭크의 각속도가 변하므로 이를 가속하고 감속하는 데 필요한 토크도 생각해야 한다. 일정한 저항력에 대하여 크랭크에 필요한 페달 토크는 다음과 같다.

\begin{equation}
\tau_p=\frac{\rho}{r}fR
-I_c\frac{r^2v_0^2}{R^2\rho^3}\rho'.
\label{eq:me-taup}
\end{equation}

첫 번째 항은 저항력을 이겨 내기 위해 필요한 토크이고, 두 번째 항은 크랭크와 체인링의 각속도가 변하는 것을 반영한 항이다. 따라서 체인링의 형상인 \(\rho(\theta)\)를 정하면 각도에 따라 필요한 페달 토크를 구할 수 있다.

이제 이 페달 토크를 만들기 위해 고관절이 얼마만큼의 토크를 내야 하는지 구하고자 한다. 크랭크의 회전각이 정해지면 페달의 위치가 정해지고, 허벅지와 정강이의 길이는 일정하므로 허벅지 각도 \(\phi\)도 정해진다. 한 번에 한 다리만 페달을 민다고 가정하였으므로 활성 다리가 반 바퀴마다 바뀐다. 이를 반영하기 위해 다리를 계산할 때의 각도는 \(\vartheta=\operatorname{mod}(\theta,\pi)\)로 두었다. 다만 체인링의 형상과 실제로 회전한 각도 \(\theta\)는 이와 같이 반 바퀴마다 접지 않는다.

허벅지의 각도를 구하기 위해 다음 세 물리량을 정의한다.

\[
\begin{aligned}
P&=2L_1(d+\ell\cos\vartheta),\\
Q&=2L_1\ell\sin\vartheta,\\
K&=L_1^2+\ell^2-L_2^2+d^2+2d\ell\cos\vartheta.
\end{aligned}
\]

허벅지와 정강이의 길이가 일정하다는 조건을 정리하면 \(P\cos\phi+Q\sin\phi=K\)이다. 이 식에서 다리의 형태에 맞는 해를 택하면 허벅지 각도와 그 변화율 \(k\)는 다음과 같다.

\begin{equation}
\phi=\arccos\left(\frac{K}{\sqrt{P^2+Q^2}}\right)+\operatorname{atan2}(Q,P),
\qquad
k=\frac{K'-P'\cos\phi-Q'\sin\phi}{Q\cos\phi-P\sin\phi}.
\label{eq:me-fourbar}
\end{equation}

여기서 프라임은 각도에 대한 미분을 뜻한다. \(\phi\)와 \(k\)를 구했으므로 크랭크가 돌 때 다리가 어떻게 움직이는지 계산할 수 있다. 다리는 질량을 가지고 있으므로 다리를 움직이는 데에도 토크가 필요하다. 이를 반영하기 위해 다리의 운동에너지가 \(\frac12J\dot{\theta}^2\)이 되도록 등가 회전관성 \(J\)를 정의하였다.

\[
\begin{aligned}
M_{\mathrm{ps}}&=\frac{L_1^2k^2+\ell^2+2L_1\ell k\cos(\phi-\vartheta)}{4},\\
u_x&=\ell\sin\vartheta-L_1\sin\phi,
&v_y&=-\ell\cos\vartheta-d+L_1\cos\phi,\\
u_x'&=\ell\cos\vartheta-L_1k\cos\phi,
&v_y'&=\ell\sin\vartheta-L_1k\sin\phi,\\
\beta'&=\frac{u_xv_y'-v_yu_x'}{L_2^2}.
\end{aligned}
\]

\begin{equation}
J=\frac{m_1L_1^2k^2}{3}
+m_2M_{\mathrm{ps}}
+\frac{m_2L_2^2(\beta')^2}{12}.
\label{eq:me-inertia}
\end{equation}

\(M_{\mathrm{ps}}\)는 정강이 질량중심의 각도에 대한 위치 변화율의 크기를 제곱한 값이다. 이 식의 첫 항은 허벅지의 회전에 의한 항이며, 두 번째와 세 번째 항은 각각 정강이의 질량중심 운동과 질량중심을 기준으로 한 회전에 의한 항이다. 크랭크의 관성 \(I_c\)는 이미 페달 토크를 구할 때 반영하였으므로 다리의 관성 \(J\)에 다시 더하지 않는다.

다리의 관성은 크랭크의 각도에 따라 변한다. 따라서 다리의 관성과 그 변화까지 생각하면 고관절이 가해야 하는 토크는 다음과 같다.

\begin{equation}
\tau_h=\frac{\tau_p+J\ddot{\theta}+\frac12J'\dot{\theta}^2}{k}.
\label{eq:me-tauh}
\end{equation}

이제 고관절 토크를 이용하여 사람의 부담을 정의하고자 한다. 같은 토크를 내더라도 그 순간 사람이 낼 수 있는 최대 토크가 작다면 더 큰 부담이 될 것이다. 따라서 그 순간의 고관절 토크 크기를 그 순간 낼 수 있는 최대 토크로 나누어 부담 \(u\)를 정의하였다. 고관절의 각속도가 커질수록 낼 수 있는 최대 토크가 직선적으로 작아진다고 가정하면 다음과 같다.

\begin{equation}
\tau_{\max}=
\max\left[\tau_{\mathrm{iso}}
\left(1-\frac{|k\dot{\theta}|}{\omega_{\max}}\right),\,0\right],
\qquad
u=\frac{|\tau_h|}{\tau_{\max}}.
\label{eq:me-burden}
\end{equation}

\(u\)가 1이면 그 순간의 토크가 가정한 최대 토크와 같다는 뜻이고, 1보다 크면 최대 토크를 넘는다는 뜻이다. 다만 이 최대 토크는 고관절의 각도나 토크의 방향에 따른 차이를 무시한 모델이며, \(u\)도 사람의 피로도를 직접 측정한 값은 아니다. 최대 토크를 구하기 위한 \(\tau_{\mathrm{iso}}\)와 \(\omega_{\max}\)는 사람을 측정하여 정한 값이 아니라 계산을 위한 가정값이다.

\subsection{최적의 체인링 찾기}

연구의 목적은 체인링의 형상을 바꾸어 사람의 부담이 한쪽 구간에 몰리지 않도록 하는 것이다. 따라서 한 바퀴 동안의 부담이 얼마나 변하는지를 목적함수로 정하고, 그 값이 작아지는 체인링을 찾고자 한다.

먼저 체인링의 형상 \(\rho(\theta)\)를 푸리에 급수로 나타낸다. 그러면 모든 각도에서의 \(\rho\)를 따로 정하는 대신 유한한 개수의 계수를 정하여 형상을 바꿀 수 있다.

\begin{equation}
\rho(\theta)=\rho_0\left[
1+\sum_{n=1}^{N_H}\left(a_n\cos n\theta+b_n\sin n\theta\right)
\right].
\label{eq:me-rho}
\end{equation}

\(\rho_0\)는 평균 유효 반지름이고, \(N_H\)는 사용하는 푸리에 성분의 최대 차수이다. 평균 유효 반지름은 고정하고 계수 \(a_n\), \(b_n\)을 바꾸어 형상에 따른 부담을 비교한다. 이때 다리의 움직임은 반 바퀴마다 반복되지만 체인링은 한 바퀴 전체에 대해 정의한다.

한편 고관절 토크 식의 분모에 있는 \(k\)가 0에 가까워지면 고관절 토크가 매우 커진다. 이 짧은 구간은 관성으로 통과하여 자전거 속력의 변화가 작다고 가정하고 부담 계산에서 제외한다. 이를 위해 미리 \(k_{\min}\)을 정하고 \(|k|\geq k_{\min}\)인 구간만 사용한다. 제외 구간에서 실제로 속력이 얼마나 변하는지는 이 가정만으로 알 수 없으므로 별도의 검증이 필요하다.

구동 구간에서의 시간 평균 부담을 \(\bar{u}\)라고 하면 우리가 최소화하고자 하는 부담의 표준편차는 다음과 같다.

\begin{equation}
E=\sqrt{\overline{(u-\bar{u})^2}}.
\label{eq:opt}
\end{equation}

위의 평균은 크랭크 각도가 아니라 시간을 기준으로 한 평균이다. 크랭크가 천천히 지나는 구간은 그만큼 오래 부담을 받으므로 더 크게 반영된다. \(E\)가 작을수록 부담의 변화가 작아져 구동 구간에서 부담이 평평해진다.

다만 부담이 매우 크더라도 일정하게만 유지되면 \(E\)가 작아질 수 있다. 따라서 평균 부담 \(\bar{u}\)가 원형 체인링일 때보다 커지지 않는 조건을 함께 둔다. 또한 실제 체인링이 될 수 있도록 \(\rho\)와 곡률반지름 \(\rho+\rho''\)에 제한을 둔다. 이 조건들은 최적화에 추가한 조건이며, 부담의 평탄함만으로 자동으로 만족되는 조건은 아니다.

\subsection{수치해석}

계산할 때에는 한 바퀴를 \(N\)개의 각도로 나눈다. 양쪽 다리가 바뀌는 경계에서는 다리의 등가 관성이 불연속일 수 있으므로 경계점을 피하여 각 구간의 가운데에서 계산한다. 다리의 형상과 관성은 체인링의 모양과는 관계가 없으므로 각도에 따른 \(k\)와 \(J\)를 먼저 구한다. \(J'\)는 각도를 조금씩 바꾸어 구한 \(J\)의 차이로 근사하며, 다리가 바뀌는 경계는 넘지 않도록 한다.

그다음 푸리에 계수를 정하면 각도마다 \(\rho\)와 \(\rho'\)가 정해진다. 이를 식 (\ref{eq:me-motion}), (\ref{eq:me-taup}), (\ref{eq:me-tauh})에 넣어 각속도, 각가속도, 페달 토크와 고관절 토크를 차례로 구한다. 고관절의 각속도에 따라 최대 토크와 순간 부담을 구하고, 제외 구간을 빼고 식 (\ref{eq:opt})의 값을 계산한다.

각도에 따라 체인링을 통과하는 시간이 다르므로 시간 평균을 계산할 때에는 각 지점의 \(\rho\)에 비례하는 가중치를 사용한다. 푸리에 계수가 모두 0이면 원형 체인링이므로 이를 기준으로 평균 부담과 표준편차를 구한다. 이후 \texttt{fmincon}을 사용하여 제약 조건을 만족하면서 \(E\)를 작게 만드는 푸리에 계수를 찾는다. 이 방법은 원형에서 시작하는 수치적 탐색이므로 찾은 결과가 모든 가능한 형상 중 반드시 최적이라는 보장은 없다.
